% §2 Background.  Delivers the objects the paper is about: tree, substitution model, similarity
% and distance, the similarity graph and its Fiedler vector, clans, thm:stdr-split.  No modelling
% hypothesis is made here: the algorithm and the sampling scheme are §3, every assumption of the
% analysis is §4.
\section{Background}\label{sec:background}

This section fixes the objects the paper works with and the algorithm it accelerates.
\Cref{sec:bg-model} sets the tree and the substitution process running on it,
\Cref{sec:bg-similarity} the similarity and distance matrices they induce on the terminal nodes
and the graph those matrices define, and \Cref{sec:bg-clans} the bipartitions a tree admits
together with the result that the Fiedler vector always returns one of them.
\Cref{sec:problem} builds the reconstruction algorithm on that result and states the problem this
paper solves.

\subsection{The generative model}\label{sec:bg-model}

\subsubsection*{Tree topology}

Let $\mathcal{T} = (V, E)$ be a rooted binary phylogenetic tree whose vertex set
$V = \mathcal{X} \cup \mathcal{H}$ splits into $m$ observed terminal nodes
$\mathcal{X} = \{x_1, \dots, x_m\}$ and $m-1$ unobserved internal nodes
$\mathcal{H} = \{h_1, \dots, h_{m-1}\}$, among them a distinguished root
$r \in \mathcal{H}$ \cite{semple2003phylogenetics}. In the phylogenetic reading the terminal
nodes are extant organisms and the internal nodes their extinct ancestors.

Every edge $e = (h, h') \in E$ is assigned a \emph{branch length} $t_e \ge 0$, a measure of the
evolutionary divergence between the two adjacent nodes it joins. We represent branch lengths
through absolute node heights: each node $h$ carries a height $\tau_h \ge 0$, increasing toward the
root, with $t_{(h,h')} = |\tau_h - \tau_{h'}|$.

\begin{definition}[Most recent common ancestor]\label{def:mrca}
For two nodes $h, h'$ of $\mathcal{T}$, the \emph{most recent common ancestor}
$\mathrm{MRCA}(h,h')$ is the ancestor of both of smallest height, equivalently the unique node of
smallest height common to the path from $h$ to the root and the path from $h'$ to the root
\cite{semple2003phylogenetics}.
\end{definition}

\noindent For a pair of terminal nodes $x_i, x_j$ we denote
\begin{equation}\label{eq:mrca_height}
   \tau_{ij} := \tau_{\mathrm{MRCA}(x_i,x_j)},
   \qquad
   d_{ij} := \sum_{e \in \mathrm{path}(x_i,x_j)} t_e ,
\end{equation}
for the height of their most recent common ancestor and for the total length of the path joining
them. Closely related terminal nodes have an MRCA of small height and a short path; the root is the
highest node in the tree. \Cref{fig:tree_topology} illustrates both conventions on a tree with $m = 7$ terminal
nodes.

The two edges adjacent to the root define the \emph{root split} of the tree: the terminal set
$\mathcal{X}$ is partitioned into the two subtrees hanging off the root. A single bipartition
$\mathcal{X} = A \cup B$, rather than the full branching pattern of $\mathcal{T}$, is what one
top-down step has to resolve; the remaining structure follows recursively.

We denote by $|A|$ and $|B|$ the sizes of the two sides, labelled so that $|A| \le |B|$ and
$|A| + |B| = m$, and we order the terminal nodes along the split,
$A = \{x_1, \dots, x_{|A|}\}$ and $B = \{x_{|A|+1}, \dots, x_m\}$. The ordering is a relabelling and
costs nothing: it is what makes each side a contiguous range of indices, so that the split appears
in the similarity matrix of \Cref{sec:bg-similarity} as a $2 \times 2$ arrangement of blocks
\eqref{eq:block_forms} rather than as a scattered set of entries.

\begin{figure}[t]
\centering
\footnotesize
\tikzset{
  inode/.style={circle, draw, fill=black!8, minimum size=6.6mm, inner sep=0pt,
                font=\tiny},
  lfN/.style={circle, draw, fill=black!8, minimum size=6.6mm, inner sep=0pt,
              font=\tiny},
  % endpoint of a highlighted path: contour in that path's colour
  lfblue/.style={lfN, draw=blue!70!black, line width=1.1pt},
  lforng/.style={lfN, draw=orange!85!black, line width=1.1pt},
  lfA/.style={circle, draw, fill=blue!22, minimum size=6.6mm, inner sep=0pt,
              font=\tiny},
  lfB/.style={circle, draw, fill=red!22, minimum size=6.6mm, inner sep=0pt,
              font=\tiny},
  br/.style={thick, black!65},
  hl/.style={line width=1.6pt, blue!70!black},
  hlr/.style={line width=1.6pt, orange!85!black},
  guide/.style={densely dotted, black!45},
  cutB/.style={line width=1pt, green!45!black, dash pattern=on 3pt off 2pt},
}

%----------------------------------------------------------------- panel (a)
\begin{subfigure}[t]{0.49\textwidth}
\centering
\begin{tikzpicture}[x=1cm, y=1.15cm]
  % leaves -- neutral fill: this panel is about heights, not clans
  \node[lfblue] (x1) at (-2.35,0) {$x_1$};
  \node[lfblue] (x2) at (-1.55,0) {$x_2$};
  \node[lforng] (x3) at (-0.75,0) {$x_3$};
  \node[lfN]    (x4) at ( 0.30,0) {$x_4$};
  \node[lforng] (x5) at ( 1.10,0) {$x_5$};
  \node[lfN] (x6) at ( 1.90,0) {$x_6$};
  \node[lfN] (x7) at ( 2.70,0) {$x_7$};
  % internal; the two MRCAs carry the colour of their path
  \node[inode, fill=blue!30]   (u)  at (-1.95,0.80) {};
  \node[inode]                 (hA) at (-1.35,1.80) {$h_A$};
  \node[inode]                 (w)  at ( 0.70,0.80) {};
  \node[inode]                 (z)  at ( 2.30,0.80) {};
  \node[inode]                 (hB) at ( 1.50,1.35) {$h_B$};
  \node[inode, fill=orange!35] (r)  at ( 0.10,2.45) {$r$};
  % branches
  \draw[br] (r)--(hA); \draw[br] (r)--(hB);
  \draw[br] (hA)--(u); \draw[br] (hA)--(x3);
  \draw[br] (u)--(x1); \draw[br] (u)--(x2);
  \draw[br] (hB)--(w); \draw[br] (hB)--(z);
  \draw[br] (w)--(x4); \draw[br] (w)--(x5);
  \draw[br] (z)--(x6); \draw[br] (z)--(x7);
  % two edge-disjoint paths: within-clan x1 -- u -- x2, cross-clan x3 -- r -- x5
  \draw[hlr] (x3)--(hA)--(r)--(hB)--(w)--(x5);
  \draw[hl]  (x1)--(u)--(x2);
  % height axis
  \draw[black!55, -{Latex[length=1.4mm]}] (-3.15,-0.05) -- (-3.15,2.85)
        node[above, font=\scriptsize, black!55] {$\tau$};
  \foreach \y/\lbl in {0/$0$, 0.80/$\tau_{12}$, 2.45/$\tau_r$}{
     \draw[black!55] (-3.25,\y) -- (-3.05,\y);
     \node[left=0pt, font=\scriptsize, black!70] at (-3.25,\y) {\lbl};
  }
  \draw[guide] (-3.05,0.80) -- (u);
  \draw[guide] (-3.05,2.45) -- (r);
  \draw[guide] (-3.05,0) -- (x1);
  % annotations
  \node[font=\scriptsize, blue!70!black] at (-1.95,-0.62) {$\mathrm{MRCA}(x_1,x_2)$};
  \node[font=\scriptsize, orange!75!black, align=left, anchor=west] at (0.45,2.82)
        {$\mathrm{MRCA}(x_3,x_5)=r$\\[-1pt] so $\tau_{35}=\tau_r$};
\end{tikzpicture}
\subcaption{Heights and coalescence}\label{fig:topology-heights}
\end{subfigure}
\hfill
%----------------------------------------------------------------- panel (b)
\begin{subfigure}[t]{0.49\textwidth}
\centering
\begin{tikzpicture}[x=1cm, y=1.15cm]
  \node[lfB] (x1) at (-2.35,0) {$x_1$};
  \node[lfB] (x2) at (-1.55,0) {$x_2$};
  \node[lfB] (x3) at (-0.75,0) {$x_3$};
  \node[lfB] (x4) at ( 0.30,0) {$x_4$};
  \node[lfB] (x5) at ( 1.10,0) {$x_5$};
  \node[lfA] (x6) at ( 1.90,0) {$x_6$};
  \node[lfA] (x7) at ( 2.70,0) {$x_7$};
  \node[inode] (u)  at (-1.95,0.80) {};
  \node[inode] (hA) at (-1.35,1.80) {};
  \node[inode] (w)  at ( 0.70,0.80) {};
  \node[inode] (z)  at ( 2.30,0.80) {};
  \node[inode] (hB) at ( 1.50,1.35) {};
  \node[inode] (r)  at ( 0.10,2.45) {$r$};
  \draw[br] (r)--(hA); \draw[br] (r)--(hB);
  \draw[br] (hA)--(u); \draw[br] (hA)--(x3);
  \draw[br] (u)--(x1); \draw[br] (u)--(x2);
  \draw[br] (hB)--(w); \draw[br] (hB)--(z);
  \draw[br] (w)--(x4); \draw[br] (w)--(x5);
  \draw[br] (z)--(x6); \draw[br] (z)--(x7);
  % one edge cut, above z: C1 = {x6,x7}, C2 = everything else
  \coordinate (m1) at ($(hB)!0.5!(z)$);
  \draw[cutB] ($(m1)!0.45cm!90:(z)$) -- ($(m1)!0.45cm!-90:(z)$);
  \node[font=\scriptsize, green!45!black, align=left, anchor=west] at (2.95,1.35)
        {one edge\\[-1pt] cut};
  \node[font=\scriptsize, green!45!black] at (2.30,-0.62) {$C_1$};
  \node[font=\scriptsize, green!45!black] at (-0.55,-0.62) {$C_2$};
\end{tikzpicture}
\subcaption{Clans are edge cuts}\label{fig:topology-clans}
\end{subfigure}

\caption{A tree with $m=7$ terminal nodes. \textbf{(a)} Node heights $\tau_h$ increase toward the
root, and $\tau_{ij}$ is the height of $\mathrm{MRCA}(i,j)$
\eqref{eq:mrca_height}. The blue path joins a pair that coalesces early; the orange path joins a
pair whose most recent common ancestor is the root, so $\tau_{35}=\tau_r$. Every pair separated by
the root behaves the same way, which is what collapses a whole block of $S$ to a single value in
\Cref{prop:cfbm}. \textbf{(b)} Removing a single edge leaves a clan on each side
(\Cref{def:clan}); the cut shown gives $C_1=\{x_6,x_7\}$ and $C_2=\{x_1,\dots,x_5\}$.}
\label{fig:tree_topology}
\end{figure}

\subsubsection*{The substitution process}

Sequence data are generated by running a Markov substitution process down the edges of
$\mathcal{T}$. The process is fixed by a single rate matrix. We work with the \emph{General
Time-Reversible} (GTR) model, the standard reversible substitution model for aligned sequence data
\cite{semple2003phylogenetics}: sequences of length $\ell$ over a $d$-character alphabet evolve
along $\mathcal{T}$ under a rate matrix $Q \in \R^{d \times d}$ satisfying
\begin{enumerate}[label=(\roman*), leftmargin=2.2em, itemsep=1pt, topsep=3pt]
\item $Q_{ij}\ge0$ for $i\ne j$, and
\item $\sum_j Q_{ij}=0$ for all $i$;
\item there is a stationary distribution $\pi>0$ with $\pi Q=0$; and
\item $\mathrm{diag}(\pi)\,Q$ is symmetric.
\end{enumerate}
Conditions (i)--(ii) make $Q$ a \emph{generator}: its off-diagonal entries are the instantaneous
rates at which one character is replaced by another, and its rows sum to zero so that
$\exp(Qt)$ is a stochastic matrix for every $t\ge0$. Condition (iii) says the character
frequencies $\pi$ are left unchanged by the process, and (iv) makes it \emph{reversible} --- a
path observed forwards in time and the same path observed backwards have the same probability,
$\pi_a\,[\exp(Qt)]_{ab}=\pi_b\,[\exp(Qt)]_{ba}$. Together (i)--(iv) imply that all eigenvalues of
$Q$ are non-positive, and hence that $\mathrm{tr}(Q)<0$.

For two \emph{adjacent} nodes $h,h'$ the substitution matrix over the branch joining them is the
matrix exponential of the generator run for the branch length,
\begin{equation}\label{eq:transition_matrix}
   P(h \mid h') = \exp\bigl(Q\,t_{(h,h')}\bigr),
   \qquad
   P(h \mid h')_{ba} = \Prob[\,h = b \mid h' = a\,].
\end{equation}
By the Markov property the transition matrix is multiplicative along a path: for nodes joined
through an intermediate node, $P(h_1 \mid h_3) = P(h_1 \mid h_2) P(h_2 \mid h_3)$. Two nodes
separated by total branch length $t$ therefore satisfy $P = \exp(Qt)$ whether or not they are
adjacent.

\subsection{Similarity, distance and the similarity graph}\label{sec:bg-similarity}

The substitution model provides a measure of similarity between every pair of nodes of the tree,
which we now define.

\begin{definition}[Pairwise similarity]\label{def:similarity}
The similarity between two nodes is the geometric mean of the determinants of the two transition
matrices between them \cite{aizenbud2023spectral},
\begin{equation}\label{eq:similarity_def}
   S(h_i, h_j) := \sqrt{\det P(h_i \mid h_j)\,\det P(h_j \mid h_i)},
   \qquad S(h_i,h_i) = 1 .
\end{equation}
\end{definition}

\noindent This quantity is a function of the tree metric alone. For a single branch of length $t$,
reversibility makes the two transition matrices similar under $\mathrm{diag}(\pi)$, so they share a
determinant, and $\det\exp(Qt) = \exp(\mathrm{tr}(Q)\,t)$ gives $S = \exp(\mathrm{tr}(Q)\,t)$.
Because $\det$ is multiplicative and $P$ is multiplicative along paths, so is the similarity: for
any two terminal nodes,
\begin{equation}\label{eq:similarity_gtr}
   S_{ij} = \exp\bigl(\mathrm{tr}(Q)\,d_{ij}\bigr),
\end{equation}
with $d_{ij}$ the path length of \eqref{eq:mrca_height}. Since $\mathrm{tr}(Q)<0$, similarity decays
exponentially in the evolutionary distance separating the pair: nodes whose common ancestor is
recent are similar, and similarity falls off geometrically as that ancestor recedes into the past.

Inverting \eqref{eq:similarity_gtr} recovers the path metric of $\mathcal{T}$ up to a positive
scale. This is the second matrix the leaves carry, the \emph{distance matrix}
\begin{equation}\label{eq:distance_from_similarity}
   \Dmat_{ij} := -\log S_{ij} = |\mathrm{tr}(Q)|\,d_{ij},
   \qquad \Dmat_{ii}=0 .
\end{equation}
Similarity and distance are entrywise transforms of one another, so either may be computed from the
same sequence data; the analysis of this paper is carried out on the similarity matrix.

\begin{definition}[Similarity matrix and similarity graph]\label{def:simgraph}
The \emph{similarity graph} $G$ is the complete weighted graph whose vertices are the $m$ terminal
nodes of $\mathcal{T}$ and whose edge $(x_i,x_j)$ carries weight $S(x_i,x_j)$. Its weighted
adjacency matrix is the \emph{similarity matrix} $S \in \R^{m \times m}$, which consists of the
pairwise similarities over the terminal nodes, $S_{ij} := S(x_i, x_j)$, and is symmetric with unit
diagonal \cite{aizenbud2023spectral}.
\end{definition}

\noindent The split is read out of $S$ through the Laplacian of that graph.

\begin{definition}[Laplacian and Fiedler vector]\label{def:fiedler}
The \emph{combinatorial Laplacian} of the similarity graph is $L := D - S$, where $D$ is diagonal
with $D_{ii} = \sum_j S_{ij}$ \cite{aizenbud2023spectral}. Write
$0 = \lambda_1 \le \lambda_2 \le \lambda_3 \le \dots \le \lambda_m$ for its eigenvalues. The
\emph{Fiedler vector} $\vv$ is a unit eigenvector belonging to the second-smallest eigenvalue
$\lambda_2$ \cite{fiedler1975property, vonluxburg2007tutorial}.
\end{definition}

\subsection{Clans and the Fiedler split}\label{sec:bg-clans}

The bipartitions of $\mathcal{X}$ that a tree admits have a standard name.

\begin{definition}[Clan; adjacent clans]\label{def:clan}
A \emph{clan} of $\mathcal{T}$ is a subtree that can be separated from the rest of $\mathcal{T}$ by
removing a single edge of the underlying unrooted tree, together with the terminal and internal
nodes it contains \cite{wilkinson2007clades, semple2003phylogenetics}. Two disjoint clans are
\emph{adjacent} if their union is again a clan \cite{aizenbud2023spectral}, equivalently if the
roots of the two subtrees are joined to a common neighbouring node.
\end{definition}

\noindent Removing one edge always leaves a clan on each side, so the clans of $\mathcal{T}$ are
exactly its edge cuts; \Cref{fig:topology-clans} shows one. The root split of \Cref{sec:bg-model} is
one of them: $A$ and $B$ are complementary adjacent clans. When we speak of the terminal nodes of a
clan we mean the elements of $\mathcal{X}$ that the separated subtree contains.

Clans, and only clans, leave a rank signature in the similarity matrix.

\begin{lemma}[Clans are the rank-one blocks, {\cite[Lem.~3.1]{jaffe2021spectral}}]\label{lem:rank-one}
Let $A$ and $B$ partition the terminal nodes of an unrooted binary tree $\mathcal{T}$, and let
$S(A,B)$ denote the $|A| \times |B|$ submatrix of cross-pair similarities. Then $S(A,B)$ is rank-one
if and only if $A$ and $B$ are clans of $\mathcal{T}$.
\end{lemma}

\noindent \Cref{lem:rank-one} is the precise sense in which the topology of $\mathcal{T}$ is
embedded in the structure of $S$: a candidate bipartition is a split of the tree exactly when the
block of similarities it cuts collapses to rank one. With the terminal nodes ordered along the
split as in \Cref{sec:bg-model}, the within-clan pairs occupy the two diagonal blocks of
$S$ and the cross-clan pairs its off-diagonal block, and \Cref{lem:rank-one} says that this
off-diagonal block is rank-one at every split of the tree and at no bipartition that is not one.

Because the similarity \eqref{eq:similarity_gtr} is a function of pairwise distances alone, $S$ is
invariant to where the root is placed on $\mathcal{T}$: the bipartitions it can express are exactly
the edge cuts of the unrooted tree.

We can now say why the similarity matrix is the object to compute. Thresholding the Fiedler vector
of the similarity graph never returns a bipartition that is not a split of the tree.

\begin{theorem}[Consistency of the Fiedler split, {\cite[Thm.~4.2]{aizenbud2023spectral}}]\label{thm:stdr-split}
Let $G$ be the similarity graph of a binary tree $\mathcal{T}$ and let $v$ be the Fiedler vector of
$G$. Partition the terminal nodes by the sign pattern of $v$,
\begin{equation}\label{eq:sign_rule}
   C_1=\{i:\;v_i\ge0\},\qquad C_2=\{j:\;v_j<0\}.
\end{equation}
Then $C_1$ and $C_2$ are adjacent clans of $\mathcal{T}$.
\end{theorem}

\noindent \Cref{thm:stdr-split} carries no hypothesis on tree shape or on the substitution
parameters: given the exact similarity matrix, the sign rule \eqref{eq:sign_rule} always cuts the
tree along a single edge. This is the result the algorithm of \Cref{sec:problem} is built on
--- it is what makes top-down spectral recovery viable at all, and it is why the recursion can merge
its two subproblems, since each side is a genuine subtree. We write $\Pisp{S}$ for the bipartition
\eqref{eq:sign_rule} of the Fiedler vector of $L=D-S$.

